\chapter{模型评价、灵敏度分析与推广}

\section{稳健性证据汇总}

本文各问题均设置了主模型和对照模型，避免结论依赖单一算法。表\ref{tab:robust_summary}汇总了主要稳健性证据。

\begin{table}[htbp]
\centering
\footnotesize
\caption{模型稳健性证据汇总}
\label{tab:robust_summary}
\begin{tabularx}{0.98\textwidth}{C{1.2cm}X X}
\toprule
问题 & 主模型 & 稳健性证据 \\
\midrule
1 & Huber稳健仿射校准 & OLS、RANSAC、Theil--Sen对照；时间块CV；bootstrap置信区间无NaN且宽度很小 \\
2 & 速度主导的复合证据变点 & 3 h、6 h、12 h速度窗口敏感性；位移趋势、持续加速度、速度跃迁三类证据互证 \\
3 & 变量专属补齐与异常识别 & 阈值3.5--5.5敏感性检查；ElasticNet、GBDT、HistGBDT模型对照；变量贡献重复划分稳定 \\
4 & 阶段归一化基线+GBDT残差 & 阶段基线、线性残差、GBDT残差、全局直接回归、分阶段直接回归消融对照 \\
5 & 五变量组合+三级速度预警 & 全部六选五组合枚举；ElasticNet与GBDT双模型排序；3 h、6 h、12 h阈值窗口敏感性 \\
\bottomrule
\end{tabularx}
\end{table}

问题三异常阈值的敏感性表现为：连续变量Hampel阈值从3.5增至5.5时，异常数量平滑下降，未出现跳变；降雨量按0.997分位数得到17个异常，与共同异常表中降雨--表面位移同步点相匹配。问题五阈值窗口检查显示，快速阶段速度中位数在3 h、6 h、12 h窗口下均接近10 mm/h，严重预警阈值稳定。

\section{题面要求覆盖与证据链闭环}

表\ref{tab:requirement_trace}给出本文的覆盖矩阵。该表用于说明：每个输出并非孤立数值，而是由对应附件、处理流程、模型检验和正文图表共同支撑。

\begin{longtable}{p{1.0cm}p{3.2cm}p{4.7cm}p{4.2cm}}
\caption{题面要求覆盖矩阵与证据链}
\label{tab:requirement_trace}\\
\toprule
问题 & 题面关键要求 & 本文模型与结果 & 证据位置 \\
\midrule
\endfirsthead
\toprule
问题 & 题面关键要求 & 本文模型与结果 & 证据位置 \\
\midrule
\endhead
\bottomrule
\endfoot
1 & 建立A、B两组位移数据校正关系 & Huber稳健仿射模型$\hat y=-0.5264+0.8817x$，避免高阶外推不稳 & 第3章模型式、表\ref{tab:q1_model_compare} \\
1 & 交叉验证并给出量化指标 & 5折时间块CV，Huber平均MAE为1.289 mm，P95AE为7.219 mm & 表\ref{tab:q1_model_compare}和摘要 \\
1 & 填写题目表1.1 & 5个待校正值为5.762、15.809、73.836、108.415、147.311 mm，并给出bootstrap区间 & 表\ref{tab:q1_result} \\
2 & 给出两个阶段转换节点 & 复合证据主节点为第8108号和第9501号，并区分阶段起点与速度跃迁点 & 表\ref{tab:q2_candidates}、表\ref{tab:q2_stage} \\
2 & 排除噪声和瞬时跳变干扰 & Hampel去尖峰、6 h滚动中位速度、持续加速度共同判别，单点尖峰不作为阶段转换 & 第4.1--4.2节、图\ref{fig:q2_segmentation} \\
2 & 各阶段建立模型、给出参数并检验 & 三段分别选取二次或三次多项式，最小$R^2$为0.99992，最大RMSE为1.422 mm & 表\ref{tab:q2_stage_model} \\
2 & 计算阶段平均速度 & 三阶段平均速度为0.262、1.859、6.234 mm/h，速度等级清晰分离 & 表\ref{tab:q2_stage} \\
3 & 数据去噪与缺失补齐 & 连续状态量采用局部线性状态空间补齐，稀疏驱动量保留真实事件强度 & 第5.1节、图\ref{fig:q3_trace} \\
3 & 单变量异常和共同异常 & 单变量异常63个，共同异常17个；共同异常按同编号至少两个变量异常定义 & 表\ref{tab:q3_single_anomaly}、表\ref{tab:q3_common_anomaly} \\
3 & 分析多源变量贡献并估计实验集表面位移 & GBDT贡献排序为孔隙水压力、深部位移、降雨量、微震事件数；实验集输出5000个估计值 & 表\ref{tab:q3_contribution}、图\ref{fig:q3_exp_scatter} \\
4 & 构建分阶段预测模型 & 阶段归一化PCHIP基线描述整体变形，GBDT残差解释扰动偏差，预测曲线单调 & 第6.1--6.2节、图\ref{fig:q4_baseline} \\
4 & 预测实验集5个指定时刻 & 预测值为5.163、30.767、60.259、383.185、383.697 mm，并给出95\%区间 & 表\ref{tab:q4_prediction}、图\ref{fig:q4_prediction} \\
5 & 从6个变量选出最优5个 & 枚举全部6种五变量组合，最终剔除干湿入渗系数，保留降雨、孔压、微震和两项爆破特征 & 表\ref{tab:q5_varselect} \\
5 & 建立以速度为指标的预警机制 & 6 h平滑速度三级阈值为2.980、6.418、9.949 mm/h，并叠加持续时间和逆速度趋势 & 表\ref{tab:q5_threshold}、图\ref{fig:q5_warning_timeline} \\
\end{longtable}

从闭环关系看，本文每个关键结论均满足三层支撑：第一层是题面附件中的原始观测；第二层是由脚本生成的中间表格、模型参数和图件；第三层是正文中的解释、对照和敏感性分析。若某个数值只出现在正文而不存在于输出表，则不纳入最终结论。该规则降低了人工转录和选择性叙述风险，也使论文能够经受复算审查。

\section{数据溯源和数值一致性审查}

为防止模型结果与附件脱节，本文对5个附件进行了逐项溯源。附件1和附件2均为10000行等间隔时序；附件3含训练集10000行、实验集5000行；附件4含训练集10000行、实验集5000行；附件5含9450行，并包含6个候选变量。所有脚本均通过统一配置读取中文附件路径，核心输出写入\texttt{code/outputs/tables}、\texttt{code/outputs/models}和\texttt{code/outputs/award}，正文只摘取这些文件中的关键数值。

\begin{table}[htbp]
\centering
\footnotesize
\caption{关键数值一致性复核}
\label{tab:numerical_audit}
\begin{tabularx}{0.98\textwidth}{C{1.0cm}X X}
\toprule
问题 & 复核项目 & 复核结论 \\
\midrule
1 & Huber校准值与附件1重算结果 & 由附件1重算得到$\beta_0=-0.5264,\ \beta_1=0.8817$，表\ref{tab:q1_result}数值一致 \\
2 & 复合节点与阶段模型来源 & 第8108号、第9501号来自增强审计输出，分阶段模型参数同步写入\texttt{q2\_final\_stage\_models.csv} \\
3 & 补齐、异常与实验集估计 & 预处理后训练集无缺失，非负变量无负值；实验集5000行均给出表面位移估计 \\
4 & 指定时刻是否真实存在 & 题目要求的5个时刻均存在于附件4实验集，预测序列单调不降且给出区间 \\
5 & 五变量组合是否枚举完整 & 六选五组合共6种，全部参与时间块交叉验证；最终阈值严格递增 \\
\bottomrule
\end{tabularx}
\end{table}

上述复核说明，本文没有使用脱离附件的人工构造序列。需要人工写入正文的数值只保留三位小数，其更高精度版本保存在CSV或JSON文件中。对于节点、异常编号和指定时刻预测值这类容易被质疑的结果，正文均同时给出表格或图形位置，保证“原始数据--算法输出--论文叙述”三者可相互追踪。

\section{科学性审查：模型必要性与可证伪性}

从科学建模角度看，一个模型是否可靠，不仅取决于误差是否小，还取决于它是否有明确适用条件，是否能被数据反驳。本文的五个主模型均满足可证伪要求。问题一若高阶模型在时间块交叉验证中明显优于Huber一次模型，则本文的线性校准会被否定；实际结果显示Huber在MAE和P95AE上更稳。问题二若第8108号之后速度不能持续抬升，则复合节点不成立；实际速度和后续平均速度均支持该节点。问题三若统一Hampel规则能在不误删降雨的情况下取得更稳定贡献排序，则变量专属处理不必要；敏感性结果说明稀疏驱动量需要单独处理。问题四若全局直接回归优于阶段归一化模型且不产生单调性问题，则阶段基线可被替代；消融实验显示阶段结构具有更好泛化解释。问题五若剔除其他变量组合显著优于当前组合，则最优变量选择会改变；全组合验证给出了明确排序。

这种可证伪设计使本文不是“先选模型再解释结果”，而是让每个模型选择都接受反事实检验。对于竞赛论文而言，这一点比单纯引入更复杂的深度模型更重要：评委不仅要看到预测值，还要看到为什么这些预测值可信、在哪些条件下可能失效，以及本文如何用敏感性分析提前暴露这些风险。

\section{模型优点}

本文模型具有以下优点。

\begin{enumerate}
  \item 可解释性强。阶段划分以速度和加速度为主，预测以阶段基线为主，预警以速度阈值为主，均能对应边坡变形机理。
  \item 抗异常能力强。Huber校准、Hampel去尖峰、状态空间补齐和稀疏事件识别分别处理不同噪声来源，避免“一刀切”清洗。
  \item 泛化证据充分。关键模型均使用时间块交叉验证，问题四还通过消融实验证明阶段归一化结构优于全局直接回归。
  \item 结果可复现。所有表格和图件由脚本生成，原始结果和增强审查结果分目录保存，便于复核。
\end{enumerate}

\section{误差来源分析}

本文模型的主要误差来源有三类。第一类是量测误差，包括位移计短时跳变、孔压传感器漂移和微震计数离散性。该误差通过Huber、Hampel和状态空间模型削弱，但无法完全消除。第二类是外部扰动未观测误差，例如地下水通道、局部开挖和坡体材料非均质性，这些因素会表现为残差模型中的不可解释波动。第三类是阶段迁移误差，即训练集与实验集虽然阶段标签一致，但各阶段持续时间和扰动强度不完全相同。问题四采用$\tau$归一化正是为了降低这类误差。

从数值结果看，问题四第三阶段预测区间明显宽于前两阶段，说明模型已经把快速阶段不确定性显式暴露出来，而不是给出过度自信的点预测。问题五严重预警也要求持续6 h和逆速度窗口条件，目的同样是控制快速阶段噪声误报。

\section{局限性}

模型也存在局限。

\begin{enumerate}
  \item 阈值具有工点依赖性。2.980、6.418和9.949 mm/h适用于本题数据，迁移到其他边坡时应重新按历史阶段速度分布标定。
  \item 爆破影响采用经验衰减特征。若能获得更精细的爆破位置、装药结构和地层参数，可引入能量衰减或波传播模型。
  \item 逆速度法用于辅助预警而非确定失稳时间。实际工程中仍需结合现场巡查、降雨预报和地质条件综合判断。
\end{enumerate}

\section{推广应用}

对于同类边坡监测任务，可按如下流程迁移本文方法：先用历史稳定期数据完成传感器校准，再以速度序列识别阶段；随后在各阶段内拟合归一化位移基线，并用降雨、孔压、微震和施工扰动解释残差；最后以阶段速度分布设定分级阈值，并用逆速度下降趋势筛查快速阶段风险。该流程只需重新估计参数，不依赖题目中特定编号，因此具有较好的工程推广性。
